\section{Round-twenty-three reconstruction: ordered marked histories and absolutely continuous entropy recovery}
\label{sec:r23-a3}

The state used in the preceding revision retained an empirical transition
measure but forgot temporal order.  It therefore could not determine the
concatenated physical path.  Its recovery also replaced a continuous
reference kernel by representative atoms, which has infinite relative
entropy.  We repair both defects by putting the order coordinate and the
physical excursion itself into the state and by coarse graining with
conditional reference densities rather than atoms.

\subsection{The exact history kernel}

Let $E$ be the Polish space of legal induced excursions.  A mark
$m\in E$ contains
\[
 m=(R(m),\tau(m),\kappa(m),\gamma_m,u_m),
\]
where $R\in\mathbb N$ is the collision length, $\tau>0$ is physical time,
$\kappa\in\mathbb Z^2$ is displacement, $u_m$ is the unstable entrance
coordinate, and
$\gamma_m:[0,\tau(m)]\to M$ is the complete billiard trajectory with its
one-sided singularity labels.  The metric on $E$ is the sum of the discrete
metrics on $(R,\kappa)$, the Euclidean metrics on $(\tau,u)$, and the uniform
metric on the time-rescaled path.  Thus convergence of marks includes
convergence of the actual physical excursion.

The complete-past space is
\[
 \mathsf H=\{h=(\ldots,m_{-2},m_{-1}):m_j\in E
             \text{ and all concatenations are legal}\}
\]
with weighted metric
$d_\sigma(h,h')=\sum_{j\ge1}\sigma^j(1\wedge d_E(m_{-j},m'_{-j}))$.
Let $U(h,m)$ append $m$ and shift the past.  Disintegration of Liouville
measure over stable leaves gives a regular conditional probability
$K(h,dm)$ supported on legal marks and the Markov kernel
\[
 \mathsf P(h,dh')=\int 1_{\{U(h,m)\in dh'\}}K(h,dm).
\]
No one-step countable surrogate replaces this continuous kernel.

Set
\[
 \chi(m)=R(m)+\tau(m)+|\kappa(m)|+
 \sup_{0\le s\le\tau(m)}d(\gamma_m(s),\gamma_m(0)).
\]
The weighted branch estimate and stable-holonomy distortion from the A2
construction imply the following nonuniform, but Lyapunov-controlled, bound.

\begin{lemma}[History-weighted exponential mark estimate]
\label{lem:r23-a3-moment}
There are $\eta>0$, $0<\lambda<1$, $C<\infty$, and a continuous
$W:\mathsf H\to[1,\infty)$ such that
\[
 \int e^{\eta\chi(m)}K(h,dm)\le C W(h),
 \qquad \mathsf P W(h)\le\lambda W(h)+C.                 \tag{A3.1}
\]
Moreover, on each set $\{W\le L\}$, the density ratio of $K(h,\cdot)$ and
$K(h',\cdot)$ on a common regular branch is bounded by
$\exp(C_Ld_\sigma(h,h'))$.
\end{lemma}

\begin{proof}
Disintegrate branchwise over the induced map.  A2's weighted branch bound
sums the inverse Jacobian multiplied by $e^{\eta(R+\tau+|\kappa|)}$ for
sufficiently small $\eta$.  The conditional density on a stable leaf differs
from the invariant branch density by the stable holonomy Jacobian.  Its
logarithm is the convergent sum of future distortion differences and is
bounded by a constant times the weighted past distortion.  Define $W$ as the
exponential of that distortion plus the weighted remote-mark cost.  Appending
one mark shifts every old contribution by the contraction factor $\sigma$;
the new contribution is controlled by the branch exponential moment.  This
gives (A3.1).  Notice that the estimate is not the false assertion
$\sup_hKe^{\eta\chi}<\infty$; histories near singularity are paid for by
$W(h)$.
\end{proof}

\subsection{The ordered stopped state}

For a legal mark sequence $(m_k)_{k\ge1}$ let
\[
 C_k=\sum_{j\le k}R(m_j),\qquad T_k=\sum_{j\le k}\tau(m_j).
\]
For collision horizon $N$ set
$\nu_N^c=\inf\{k:C_k\ge N\}$, and for physical horizon $N$ set
$\nu_N^t=\inf\{k:T_k\ge N\}$.  We describe the collision version; the
physical version is identical after replacing $C$ by $T$.

The ordered marked point measure is
\[
 \mathcal O_N^c={1\over N}
 \sum_{k=1}^{\nu_N^c}
 \delta_{(C_{k-1}/N,m_k)}
 \quad\text{on }[0,1]\times E.                           \tag{A3.2}
\]
The first coordinate records order.  The stopped physical path
$\Gamma_N^c\in D([0,1];M)$ is obtained by concatenating the $\gamma_{m_k}$ in
that order and truncating the final excursion at collision time $N$.  The
terminal object $\zeta_N^c$ is the full final mark together with the truncation
fraction; it is not a zero-probability point event.  Finally, the recession
measure
\[
 \mathcal R_N^c={1\over N}
 \sum_{k\le\nu_N^c}\chi(m_k)
 1_{\{\chi(m_k)>K_N\}}
 \delta_{(C_{k-1}/N,[m_k])},
 \qquad K_N\uparrow\infty,\quad K_N/N\downarrow0,         \tag{A3.3}
\]
lives on the compactification of directions of legal marks.  It records
one-big-excursion trajectories rather than discarding them.

Set
\[
 Z_N^c=(\mathcal O_N^c,\Gamma_N^c,
        C_{\nu_N^c}/N,T_{\nu_N^c}/N,
        \mathcal R_N^c,\zeta_N^c).                        \tag{A3.4}
\]
The topology is weak convergence with the $\chi$-moment on the ordered
measure, Skorokhod convergence of $\Gamma$, ordinary convergence of the two
clocks, vague convergence on the recession compactification, and the pointed
mark topology for $\zeta$.  Formula (A3.2), not an unordered edge occupation,
is the primitive state coordinate.

\begin{proposition}[Ordered reconstruction]
\label{prop:r23-a3-order}
On every set on which the total $\chi$-moment is bounded and simultaneous
singular endpoints are absent, the concatenation map
$(\mathcal O,\zeta)\mapsto\Gamma$ is well defined and continuous.  For the
full state (A3.4) it is the coordinate projection.  Two distinct Eulerian
orderings with the same unordered transition counts remain distinct because
their first coordinates in (A3.2) differ.
\end{proposition}

\begin{proof}
Order the atoms by their first coordinate and concatenate their path marks.
Moment boundedness makes the contribution of marks beyond a finite compact
subset uniformly small after rescaling.  On that compact subset, endpoint and
path matching are closed legal relations and uniform convergence of each
mark gives Skorokhod convergence of the concatenation.  Ties correspond to
simultaneous endpoints and are excluded on the stated regular set; their
closure is retained separately in the singular/recession coordinate.  The
last statement is immediate from the time coordinate in (A3.2).
\end{proof}

\subsection{Predictable controls and exact entropy}

A control is a predictable kernel
$q_k(dm)=q(h_{k-1},k/N,dm)$ satisfying
$q(h,s,\cdot)\ll K(h,\cdot)$.  Its stopped cost is
\[
 \mathcal A_N(q)=\frac1N
 \mathbb E^q\sum_{k=1}^{\nu_N^c}
 H(q(h_{k-1},k/N,\cdot)\mid K(h_{k-1},\cdot)).            \tag{A3.5}
\]
The terminal truncation has no artificial point cost: it is part of the last
controlled transition and its conditional law is retained.

For a finite measurable partition $\mathcal P=\{C_1,\ldots,C_J\}$ of $E$
whose cell boundaries have zero $K(h,\cdot)$ mass, define the reference-cell
recovery
\[
 q^{\mathcal P}(h,s,dm)=
 \sum_{j:q(h,s,C_j)>0}q(h,s,C_j)
 {1_{C_j}(m)K(h,dm)\over K(h,C_j)}.                       \tag{A3.6}
\]
Every recovered mark is an actual mark in the support of the physical kernel.
There is no representative Dirac mass.

\begin{lemma}[Finite-cell recovery at finite relative entropy]
\label{lem:r23-a3-recovery}
For every control of finite cost there are refining partitions
$\mathcal P_n$ and finite-memory approximations of $h\mapsto q$ such that
$q^{\mathcal P_n}\Rightarrow q$, the induced ordered states converge, and
\[
 H(q^{\mathcal P_n}(h,s,\cdot)\mid K(h,\cdot))
 =\sum_jq(C_j)\log{q(C_j)\over K(C_j)}
 \le H(q(h,s,\cdot)\mid K(h,\cdot)).                     \tag{A3.7}
\]
The costs converge after a final martingale refinement.
\end{lemma}

\begin{proof}
Equation (A3.7) is the chain rule for relative entropy, or equivalently data
processing under the cell map.  Conditional reference measures in (A3.6) are
absolutely continuous with respect to $K$, even when $K$ is non-atomic.
Choose compact sets from (A3.1), partitions of vanishing diameter on them, and
one tail cell.  Conditional expectations of the density $dq/dK$ on the
resulting sigma-fields converge in $L^1(q)$ and their entropies increase to
$H(q\mid K)$.  Approximate the predictable density by conditional
expectations on finite past cylinders.  Since all samples remain in the legal
support of $K$, endpoint compatibility is exact.  Proposition
\ref{prop:r23-a3-order} and the moment bound give convergence of the ordered
state and path.
\end{proof}

\subsection{Laplace principle including long excursions}

For a candidate controlled law $Q$ on
$[0,1]\times\mathsf H\times E$, write $Q(ds,dh,dm)=
\ell(ds,dh)q_{s,h}(dm)$ and impose the weak history balance relation obtained
from $h\mapsto U(h,m)$.  Let $\mathfrak Z(Q)$ be the ordered state constructed
from its chronological disintegration, including its recession part.  Define
\[
 I_c(z)=\inf\left\{
 \int H(q_{s,h}\mid K(h,\cdot))\,\ell(ds,dh):
 \mathfrak Z(Q)=z\right\}.                               \tag{A3.8}
\]
There is no condition $\nu_N^c/N\ge c_0$.  A control which chooses one mark of
length $N$ has order-$N$ entropy and is represented by a nonzero recession
coordinate with the corresponding finite cost.

\begin{theorem}[Ordered collision- and physical-time LDP]
\label{thm:r23-a3-ldp}
The variables $Z_N^c$ satisfy a good Laplace principle at speed $N$ with rate
$I_c$.  The physical-time states $Z_N^t$ satisfy the analogous principle with
rate $I_t$.  Both rates include recurrent, terminal, and recession costs and
contract to the actual ordered path.  For every bounded continuous $F$,
\[
 \lim_{N\to\infty}-{1\over N}
 \log\mathbb E\exp\{-NF(Z_N^c)\}
 =\inf_z\{F(z)+I_c(z)\}.                                  \tag{A3.9}
\]
\end{theorem}

\begin{proof}
Apply the entropy variational identity successively to the exact conditional
kernels $K(h_{k-1},dm_k)$ up to the stopping index.  This gives (A3.5) for
arbitrary predictable controls, not merely stationary current-state kernels.
Bounded cost and (A3.1), combined with the entropy inequality
$q(A)\le(H(q\mid K)+\log2)/\log(1+1/K(A))$, give compact containment of the
ordered measures and histories.  Marks of size comparable with $N$ are not
forced into a compact recurrent set; their normalized mass is tight on the
recession compactification (A3.3).  The terminal conditional law is tight by
the same estimate.

For the upper bound, extract a controlled limit, pass to the history balance
relation, and use lower semicontinuity of relative entropy.  For the lower
bound, first use Lemma~\ref{lem:r23-a3-recovery}, then concatenate finite
memory cell controls.  All transitions are sampled from conditional physical
kernels, so the constructed path is legal and has finite entropy.  The
recurrent and recession pieces are joined by connector cells of total
$\chi$-mass $o(N)$; their entropy is $o(N)$ by (A3.1).  This proves (A3.9).
Goodness follows from the same compact-containment argument.  Replacing the
collision clock by the physical clock changes only the stopping functional;
the joint two-clock state makes the contraction continuous off the terminal
point, which is already included in $\zeta$.
\end{proof}

The proof is an exact-kernel entropy argument in the spirit of weak-convergence
large deviations \cite{DupuisEllis1997,FengKurtz2006}, but the legal support,
order coordinate, random stopping, and recession mark are verified here for
the billiard kernel.

\subsection{Local conditioning without a fixed-cylinder substitution}

Let $Y_N=(S_{\nu_N}X,S_{\nu_N}\tau)$ be the finite-dimensional lattice--roof
constraint from A2.  For a regular central sequence $y_N$, A2 gives a local
denominator of order $N^{-2}$.  To treat a path functional $F$, choose
finite-memory uniformly continuous approximants $F_M$ on rate sublevels.
The ordered-state exponential tightness gives
\[
 \lim_{M\to\infty}\limsup_{N\to\infty}{1\over N}
 \log\mathbb P(|F-F_M|>\delta)=-\infty.                  \tag{A3.10}
\]
For fixed $M$, insert the finite-memory Feynman--Kac weight into the A2
operator proof and repeat the local Fourier inversion; the leading local
prefactor is continuous in the finite source.  Only after this step do we let
$M\to\infty$ using (A3.10).

\begin{corollary}[Conditioned ordered path principle]
\label{cor:r23-a3-conditioned}
Under a regular A2 local constraint $Y_N=y_N$, the conditional laws of
$Z_N^c$ satisfy the Laplace principle with rate
$I_c$ minus its minimum on the constraint fibre.  The numerator is not treated
as a fixed bounded cylinder insertion; it is obtained by the two-stage
finite-memory/exponential-approximation argument above.
\end{corollary}

\begin{proof}
For each $M$, divide the source-inserted local expansion by the A2 denominator.
The polynomial $N^{-2}$ factors cancel and the exponential part is the
constrained variational problem.  Equation (A3.10) then transfers the result
to $F$.  Exposed controls are dense by Lemma~\ref{lem:r23-a3-recovery}, and
standard upper/lower regularization completes the Laplace principle.
\end{proof}

\subsection{Round-Twenty-Two regression}

The mark sequences $AB$ and $BA$ have different ordered measures (A3.2), even
when their unordered occupations and clock totals agree.  For every non-atomic
$K(h,\cdot)$, recovery (A3.6) remains absolutely continuous and has the finite
entropy (A3.7); no atom is introduced.  A controller conditioning the first
mark on $R\ge N$ is allowed, has order-$N$ cost, may satisfy
$\nu_N/N\to0$, and appears in the recession sector rather than contradicting a
false lower-tightness assertion.  Finally, path-level conditioning is proved
by (A3.10) and source-inserted local inversion, not by calling $e^{-NF}$ a
bounded cylinder observable.
